% Bagean: sajarah
% Bab: biografi
% Pérangan: emmy-noether

\documentclass[../../../include/open-logic-section]{subfiles}

\begin{document}

\olfileid{his}{bio}{noe}

\olsection{Emmy Noether}

\olphoto{noether-emmy}{Emmy Noether}

Emmy Noether (\textsc{ner}-ter) lair ing Erlangen,
Jerman, tanggal 23 Maret 1882, ing kulawarga
cendekiawan kelas menengah ndhuwur. Dipuji
minangka ``ibu aljabar modhèren,'' Noether
mènèhi sumbangan kang ngrintis dalan anyar
marang matématika lan fisika, senajan ngadhepi
pepalang gedhé tumrap pendhidhikan wanita.
Ing Jerman wektu iku, bocah wadon dikarepaké
sinau seni lan ora diidinaké mlebu sekolah
persiapan universitas. Nanging, sawisé
melu ngrungokaké kuliah ing Universitas
G\"{o}ttingen lan Erlangen (papan bapaké
dadi profèsor matématika), Noether
pungkasané bisa ndhaptar dadi mahasiswa
ing Erlangen ing taun 1904, nalika aturané
dianyari supaya mahasiswa wanita
diidinaké. Dhèwèké nggayuh gelar doktor
matématika ing taun 1907.

Senajan mumpuni, Noether ngadhepi akèh
panentangan sajroning kariré. Saka taun
1908--1915, dhèwèké mulang ing Erlangen
tanpa bayaran. Ing wektu iku, dhèwèké
narik kawigatèné David Hilbert, salah
siji matématikawan utama ing donya
wektu iku, kang ngundang dhèwèké
menyang G\"{o}ttingen. Nanging wanita
dilarang dadi profèsor, lan dhèwèké
namung bisa menehi kuliah nganggo
jeneng Hilbert, maneh-maneh tanpa
bayaran. Ing wektu iki dhèwèké
mbuktèkaké apa kang saiki diarani
teorema Noether, kang nganti saiki
isih digunakaké ing fisika téoritis.
Noether pungkasané diwènèhi hak
mulang ing taun 1919. Wangsulan
Hilbert marang panentangan kang
terus saka kolega-kolega ing
universitasé, miturut crita,
yaiku: ``Para Bapak, senat fakultas
iku dudu papan adus.''

Ing pungkasan dasawarsa 1920-an,
dhèwèké musataké pakaryané ing
aljabar abstrak, lan sumbangané
ngowahi bidhang iku kanthi dhasar.
Ing bukti-buktiné, dhèwèké kerep
nggunakaké kang diarani syarat
runtutan munggah, kang nyatakaké
yèn ora ana runtutan tanpa wates
kang mundhak kanthi ketat saka
himpunan-himpunan tartamtu. Upamané,
struktur aljabar tartamtu kang
saiki diarani gelang Noetherian
nduwèni sipat yèn ora ana
runtutan ideal tanpa wates
$I_1 \subsetneq I_2 \subsetneq \dots$.
Syarat iki bisa digeneralisasi
menyang urutan parsial apa waé
(ing aljabar, iki ngenani kasus
khusus ideal kang diurutaké
miturut relasi subhimpunan), lan
kita uga bisa nimbang syarat
runtutan mudhun kang dadi dualé:
saben runtutan ing urutan parsial
kang \emph{mudhun} kanthi ketat
pungkasané mandheg. Yèn sawijining
urutan parsial nyukupi syarat
runtutan mudhun, induksi bisa
digunakaké ing sadawané urutan
iku kaya induksi ing urutan
$<$ ing~$\Nat$. Urutan kaya
iki diarani \emph{mapan kanthi becik}
utawa \emph{Noetherian}, lan
prinsip bukti kang cocog diarani
\emph{induksi Noetherian}.

Noether iku wong Yahudi, lan
nalika Nazi ngrebut kuwasa ing
taun 1933, dhèwèké dipecat
saka jabatané. Untungé,
Noether bisa pindhah menyang
Amérika Sarékat kanggo
jabatan sauntara ing Bryn Mawr,
Pennsylvania. Nalika ana ing
kono, dhèwèké uga mènèhi
kuliah ing Princeton, sanajan
miturut panemuné universitas
iku ora ramah marang wanita
\citep[81]{Dick1981}. Ing taun
1935, Noether ngalami operasi
kanggo ngangkat tumor rahim.
Dhèwèké séda amarga infèksi
minangka akibat saka operasi
iku, lan dikubur ing Bryn Mawr.

\begin{reading}
Kanggo biografi Noether, delengen
\citet{Dick1981}. Perimeter
Institute for Theoretical Physics
nyedhiyakake kuliah-kuliahé
ngenani urip lan pangaruh
Noether kanthi daring
\citep{Perimeter2015}.
Yèn panjenengan wis kesel maca,
\emph{Stuff You Missed in History Class}
nduwèni podcast ngenani urip
lan pangaruh Noether
\citep{Frey2015}. Karya-karya
Noether kang diklumpukaké
kasedhiya ing basa Jerman
asliné \citep{Noether1983}.
\end{reading}

\end{document}
